
% =====================================================================
\section{Measurement as reshuffling, and the unreachability of identity}
\label{sec:measurement}
% =====================================================================

We now give the operational reading of the invariant residual and the monotone
record: what it is to \emph{measure} a part, and why measurement reduces ignorance
about a thing without ever delivering the thing. A measurement is a reshuffling of
the knowledge-graph that conserves the medium; it tightens the boundary of a part
and shrinks the part's interior residue toward, but never to, the floor. We prove
that resolving a part's identity to a point would require exhausting the medium,
which reshuffling conserves by construction---so identity is structurally
unreachable. The section closes with the two-pronged way this prediction is borne
out wherever finitely many parts are identified by instruments.

\subsection{Measurement reshuffles; the medium is conserved}

\begin{definition}[Reshuffling]
\label{def:reshuffle}
A \emph{reshuffling} of a contact graph $\Gph$ relative to a medium $\med$ is a
weighted isomorphism of $\Gph$ that fixes the medium: a relabelling and
re-edging of the parts that preserves all edge weights and leaves the medium's
total contact unchanged. A \emph{measurement} of a part $A$ is a reshuffling that
adds boundary edges incident to $A$---new distinctions of $A$ from the rest.
\end{definition}

\begin{proposition}[Measurement conserves the invariant]
\label{prop:measure-conserves}
A measurement, being a reshuffling \textup{(}weighted isomorphism fixing the
medium\textup{)}, leaves the residual invariant: $\Res(\Gph')=\Res(\Gph)$
\textup{(}\Cref{thm:residual-invariant}\textup{)}, and leaves the medium's
conserved content unchanged. Measurement builds new boundary from existing
material; it creates no new invariant content.
\end{proposition}

\begin{proof}
By \Cref{def:reshuffle} a measurement is a weighted isomorphism; by
\Cref{thm:residual-invariant} the residual is preserved under weighted
isomorphism, so $\Res(\Gph')=\Res(\Gph)$. The medium is fixed by the reshuffling
(\Cref{def:reshuffle}), so its total contact is unchanged. The new edges incident
to $A$ are re-arrangements within the fixed total; no invariant content is added.
\end{proof}

\subsection{Identity is the interior residue, not the boundary}

We distinguish the two faces of a part's minimum boundary. The boundary edges are
what a measurement \emph{builds} (the marks of what the part is not); the interior
residue is what a measurement \emph{cannot reach} (the part's own conserved
content). Identity is the latter.

\begin{definition}[Interior residue; boundary]
\label{def:interior}
Let $A$ be a part with minimum boundary $\cut(A)$ to the medium. Its
\emph{boundary} is the edge set $\cut(A)$---the distinctions of $A$ from the rest
\textup{(}the ``not-$A$'' relations\textup{)}. Its \emph{interior residue} is the
conserved content on $A$'s own side of the boundary that no reshuffling
transfers across $\cut(A)$: the part's private invariant
\textup{(}\Cref{def:private}\textup{)} regarded as the water $A$ holds.
\end{definition}

\begin{theorem}[Identity is the conserved interior, approached only from outside]
\label{thm:identity-interior}
For a part $A$ in a medium $\med$:
\begin{enumerate}[label=\textup{(\roman*)},leftmargin=2.2em]
\item a measurement adds edges to the boundary $\cut(A)$, tightening the
distinction of $A$ from the rest, and thereby reduces the \emph{exterior}
ambiguity of $A$ \textup{(}what $A$ might be confused with\textup{)};
\item it does not alter the interior residue, which is conserved under all
reshuffling \textup{(}\Cref{prop:measure-conserves}\textup{)};
\item hence measurement reduces how much is unsettled \emph{about the placement of
$A$ among the rest}, never delivering the interior residue itself. Identity---the
interior---is approached from outside by tightening the boundary and is never
reached.
\end{enumerate}
\end{theorem}

\begin{proof}
(i) A measurement adds boundary edges (\Cref{def:reshuffle}); each new edge
distinguishes $A$ from a further part, reducing the set of parts from which $A$ is
not yet distinguished---its exterior ambiguity. (ii) The interior residue is the
private invariant of $A$ (\Cref{def:interior,def:private}), invariant under
weighted isomorphism (\Cref{thm:private}\,(ii)), hence unchanged by the reshuffling
that the measurement is (\Cref{prop:measure-conserves}). (iii) Combining: each
measurement moves only the exterior (the boundary), never the interior; the
interior is conserved. So the sequence of measurements drives exterior ambiguity
downward while the interior stands fixed, and the interior is never produced by
the process---only ever bounded more tightly from outside.
\end{proof}

\subsection{The self-defeat of identity-resolution}

\begin{theorem}[Identity-resolution would exhaust the medium]
\label{thm:self-defeat}
Let $\med$ be a non-completable medium \textup{(}\Cref{def:noncompletable}\textup{)}
and $A$ a thing in it. To resolve $A$'s identity to a point---to drive its
interior residue to zero, $\idres(A)=0$---would require the comparison of $A$
against $\med\setminus A$ to be complete, hence the medium to be exhausted. But a
measurement is a reshuffling that conserves the medium
\textup{(}\Cref{prop:measure-conserves}\textup{)}. Therefore no measurement, and
no sequence of measurements, resolves $A$'s identity: the operation available for
the task conserves exactly the quantity the task would have to eliminate.
\end{theorem}

\begin{proof}
By \Cref{thm:floor-infinitude}, $\idres(A)=0$ requires the comparison of $A$
against $\med\setminus A$ to be complete, which exhausts the medium and
contradicts non-completability. Independently, the only operation in hand---
measurement---is by \Cref{prop:measure-conserves} a reshuffling that leaves the
medium's conserved content unchanged; it cannot exhaust the medium, since it
conserves it. Thus the goal ($\idres(A)=0$) demands eliminating the medium while
the means (measurement) conserves it: the procedure is self-defeating, and
$\idres(A)>0$ persists under every measurement.
\end{proof}

\begin{corollary}[To know the thing exactly is to know the whole exactly]
\label{cor:know-whole}
Resolving $A$'s identity to a point would require knowing $\med\setminus A$
exactly---everything $A$ is compared against. Since $\med$ is non-completable and
measurement conserves it, this is unreachable. One cannot know a thing exactly
without exhausting the inexhaustible reference in which it is a thing.
\end{corollary}

\begin{proof}
Immediate from \Cref{thm:self-defeat}: completing the comparison is exhausting
$\med\setminus A$, i.e. knowing the medium exactly, which non-completability and
conservation jointly forbid.
\end{proof}

\subsection{Resolution and multiplicity: a standing falsification of point-truth}

The preceding theorems make a prediction about any system in which finitely many
parts are identified by repeated measurement: identification never terminates, and
it never terminates in two distinct ways---by depth and by direction. We record
the prediction and note that it is what is observed.

\begin{proposition}[Resolution never saturates]
\label{prop:resolution}
If the identity of a part were a point \textup{(}$\idres(A)=0$ attainable\textup{)},
there would be a finite measurement past which further measurement is inert---the
point being already reached. Under \Cref{thm:self-defeat}, $\idres(A)>0$ after
every measurement, so further measurement always reduces residual and never
becomes inert: \emph{resolution} \textup{(}the depth of distinction\textup{)} can
always be increased and never saturates.
\end{proposition}

\begin{proof}
Point-identity would make some finite measurement terminal, as resolving a point
admits no improvement past attainment. \Cref{thm:self-defeat} forbids attainment:
the interior residue is positive after every measurement, so each further
measurement strictly tightens the boundary (\Cref{thm:identity-interior}\,(i)) and
none is terminal. Hence resolution increases without saturation.
\end{proof}

\begin{proposition}[Multiplicity never terminates]
\label{prop:multiplicity}
If the identity of a part were a point, a single line of measurement would exhaust
it, and a second line of measurement on an independent basis would be redundant.
Under \Cref{thm:self-defeat}, each independent line of measurement is a distinct
family of boundary edges---a distinct complement against which $A$ is
individuated---and each reduces a distinct portion of exterior ambiguity. Hence
independent lines of measurement are never redundant: \emph{multiplicity}
\textup{(}the number of independent bases of distinction\textup{)} can always be
extended.
\end{proposition}

\begin{proof}
A point identity is exhausted by any complete comparison, so a second independent
comparison adds nothing---redundancy. By \Cref{thm:self-defeat} no comparison is
complete; by \Cref{thm:identity-interior} each line of measurement contributes its
own boundary edges against its own complement, reducing a portion of exterior
ambiguity not addressed by the others. Therefore each independent line contributes
non-redundantly, and new lines can always be added.
\end{proof}

\begin{theorem}[Two-pronged falsification of point-truth]
\label{thm:falsification}
In any system of finitely many parts identified by measurement within a
non-completable medium, point-identity is refuted in two independent ways:
resolution does not saturate \textup{(}\Cref{prop:resolution}\textup{)} and
independent lines of measurement do not become redundant
\textup{(}\Cref{prop:multiplicity}\textup{)}. Both are consequences of
$\idres(A)>0$ \textup{(}\Cref{thm:self-defeat}\textup{)}, and either alone is
inconsistent with point-identity. The identity of a part is therefore a region
\textup{(}a cell of positive interior residue\textup{)}, never a point.
\end{theorem}

\begin{proof}
Point-identity implies a terminal resolution (contradicted by
\Cref{prop:resolution}) and the redundancy of a second independent line
(contradicted by \Cref{prop:multiplicity}). Each contradiction alone refutes
point-identity; both follow from the positive interior residue of
\Cref{thm:self-defeat}. Positive interior residue is exactly cell-identity
(\Cref{thm:residual-region}). Hence identity is a cell.
\end{proof}

\begin{remark}[Interpretation: what an instrument does, on which nothing depends]
\label{rem:instrument-reading}
A reader may recognise \Cref{thm:falsification} in the practice of identifying a
substance by instruments. Were the identity of a substance a point---a single
determinate value---one instrument would suffice and no finer one would be wanted;
yet substances are identified by many instruments on incommensurable principles,
each refined across generations, none terminal. On the present reading each
instrument is a line of measurement contributing boundary edges against its own
complement; ``the substance'' is the cell these lines jointly bracket---the
conserved interior no instrument delivers. A measurement does not make the
substance more itself, nor its value more true; the invariant is conserved
(\Cref{prop:measure-conserves}). It reduces how much is unsettled about the
substance's placement among everything it is not. That resolution never saturates
and instruments never stop multiplying is, on this reading, the standing evidence
that identity is a cell and not a point: a point would need neither finer
resolution nor a second instrument. We assert only the graph theorems
(\Cref{thm:identity-interior,thm:self-defeat,thm:falsification}); the reading is
orientation and is not relied upon.
\end{remark}
